\section{Original versus Tree Collaboration}
\label{sec:tree}
\subsection{Coarse Attention and Maximum Spanning Tree}
The second custom experiment asks whether restricting expensive
fine-resolution collaboration reduces execution time and how much detection
quality is lost. The released weights are reused without retraining
\cite{recordtree}. The variant has three stages:
\begin{enumerate}
  \item \textbf{Coarse scoring.} Each $32\times32$ feature map is pooled to
    $4\times4$. All ordered pairs are aligned and scored at this coarse
    resolution, with weights normalized over senders. Each directional
    attention map is reduced to its mean squared attention.
  \item \textbf{Tree selection.} The two directional scores of an unordered
    vehicle pair are averaged. Descending Kruskal selection keeps the
    highest-scoring links that do not create a cycle, until $n-1$ remain.
    This is a maximum spanning tree, rather than a minimum spanning tree.
  \item \textbf{Fine fusion.} Each selected edge is used in both directions,
    and every self pair is retained. The original $32\times32$ alignment and
    scoring are applied only to these pairs. Attention is renormalized over
    the retained senders, and the normal decoder and heads follow.
\end{enumerate}

For each receiver--sender pair, the coarse attention statistic can be written
as
\[
 q_{u\to v}=\frac{1}{H_cW_c}\sum_{r,c}a^{\mathrm{coarse}}_{u\to v}(r,c)^2,
 \qquad q_{\{u,v\}}=\frac{q_{u\to v}+q_{v\to u}}{2},
 \quad H_c=W_c=4.
\]
The undirected tree determines which directed fine interactions are retained.
Fine pair evaluations change from $n^2$ to
\[
 n+2(n-1)=3n-2.
\]
At $n=15$, that is 225 original fine pairs versus 43 tree fine pairs, in
addition to the tree variant's 225 coarse pairs. Only the fine-pair count
becomes linear. Coarse scoring remains quadratic, and CPU Kruskal selection
introduces additional work. This is not a claim of linear complexity for the
complete algorithm.

There is one fusion round. A receiver can use its immediate tree neighbors,
but a path in the spanning tree does not transmit every remote vehicle's
features across multiple hops during that round. Sparsity therefore changes
the available information as well as the amount of computation.

\begin{figure}[htbp]
\centering
\begingroup
\colorlet{MoverCyan}{ReportCyan}
\begin{tikzpicture}[font=\small]
  \node[font=\bfseries] at (2,4.0) {Original: complete graph};
  \node[font=\bfseries] at (10.2,4.0) {Selected tree};
  \foreach \Agent in {1,...,5} {
    \pgfmathsetmacro{\Angle}{90+360*(\Agent-1)/5}
    \coordinate (full-node-\Agent) at
      ({2+1.45*cos(\Angle)},{2.15+1.45*sin(\Angle)});
    \coordinate (tree-node-\Agent) at
      ({10.2+1.45*cos(\Angle)},{2.15+1.45*sin(\Angle)});
  }
  \foreach \Sender in {1,...,5} {
    \foreach \Receiver in {1,...,5} {
      \ifnum\Sender<\Receiver
        \draw[{Stealth[length=1.2mm]}-{Stealth[length=1.2mm]},
          draw=ReportNavy!40,line width=0.6pt,
          shorten <=8pt,shorten >=8pt]
          (full-node-\Sender) -- (full-node-\Receiver);
      \fi
    }
  }
  % The saved n=5 preflight links from the versioned topology-example.tex.
  \foreach \From/\To in {3/4,1/2,2/3,4/5} {
    \draw[{Stealth[length=1.2mm]}-{Stealth[length=1.2mm]},
      draw=ReportCyan,line width=1pt,shorten <=8pt,shorten >=8pt]
      (tree-node-\From) -- (tree-node-\To);
  }
  \foreach \Agent in {1,...,5} {
    \foreach \Prefix in {full,tree} {
      \node[circle,fill=ReportNavy,draw=ReportCyan,line width=0.8pt,
        minimum size=0.45cm,inner sep=0pt,text=white,font=\bfseries]
        at (\Prefix-node-\Agent) {\Agent};
    }
  }
  \draw[-{Stealth[length=2mm]},draw=ReportNavy,line width=0.9pt]
    (4.0,2.15) -- (8.2,2.15);
  \node[align=center] at (6.1,1.65) {Coarse scoring\\and Kruskal selection};
  \node at (2,0.35) {$n^2=25$ fine pairs};
  \node at (10.2,0.35) {$3n-2=13$ fine pairs};
\end{tikzpicture}
\endgroup

\caption[Complete collaboration and the saved five-slot tree]{The project
topology diagram adapted for the report, using the saved five-slot preflight
tree. Double arrows denote interactions in both directions. Self links
are counted in the computation but omitted from the drawing. The original
has 25 fine pairs; the tree has 13 fine pairs plus coarse scoring.}
\label{fig:tree-topology}
\end{figure}

Figure~\ref{fig:tree-topology} gives a concrete example of what a receiver
loses when edges are removed. In the displayed tree, node 1 has node 2 as
its only peer neighbor. Its fine fusion can combine the self contribution
and node 2, while the original graph can use all five inputs. Node 3 has
two peer neighbors and can combine those with its own features. A path
from node 1 to node 5 establishes graph connectivity, but it does not
make node 5's original map an immediate input to node 1 in one fusion round.

The displayed topology is one saved preflight example. Selection is driven
by the coarse attention scores of the evaluated input group; the figure
does not assign a fixed chain to every frame. It illustrates the difference
between retaining connectivity and retaining all direct fine observations.

\subsection{Matched Timing and Accuracy Protocol}
The comparison uses the same 30 ego-1 views from three recorded test scenes
at every condition, with 427 reference car boxes counted once across the
cohort. Both methods use the unchanged checkpoint, float32 precision, and
the same retained-prediction settings. There are three timing trials and 90
samples per method and $n$, covering $n=1,\ldots,15$. The specification records
twenty warmups per method/count/trial, shuffled condition order with seed
zero, deterministic cuDNN, and TF32 disabled.

The timer includes the complete model, coarse scoring, CPU tree selection,
fine fusion, decoder, and heads. Inputs are already on the GPU. Reads and
packing, transfers, warmups, model loading, box decoding, NMS, and AP scoring
are outside the forward timer. AP is evaluated separately on the original
ego-1 labels. The run contains 2,700 timed forwards in total, and recorded ego
outputs are identical across timing trials for a given condition.

As in the runtime-only workload, $n\leq5$ uses distinct recorded inputs,
while larger $n$ replays the five recorded scans and poses. AP above five
describes the behavior of those repeated inputs, not the benefit of extra
independently recorded viewpoints.

\subsection{Complete Results and Trade-Off}
Table~\ref{tab:tree} reproduces all fifteen measured conditions from the
versioned CSV and checks them against the saved JSON during report builds.
The AP values are percentages; differences between them are percentage points.

\begin{table}[htbp]
\centering
\caption{Matched Original/Tree comparison on 30 ego-1 views. Below the rule
after $n=5$, additional slots replay the five recorded inputs. The table is
separate from the full-test benchmark and the runtime-only sweep.}
\label{tab:tree}
\input{tree-table.generated}
\end{table}

At $n=15$, median time falls from \ReportOriginalTime\ to \ReportTreeTime\ ms,
a \ReportTreeReduction\% reduction. AP@0.5 decreases from
\ReportOriginalAPFifty\% to \ReportTreeAPFifty\%, a loss of
\ReportTreeDropFifty\ percentage points. AP@0.7 decreases from
\ReportOriginalAPSeventy\% to \ReportTreeAPSeventy\%, a loss of
\ReportTreeDropSeventy\ points. This is a substantial runtime reduction
accompanied by a substantial accuracy loss, rather than an improvement at
unchanged detection quality.

The same trade-off is visible within the distinct-input region. At $n=5$,
median time falls from 41.57 to 36.54 ms, while AP@0.5 falls from 74.46\% to
62.65\% and AP@0.7 from 67.53\% to 53.66\%. Thus the accuracy penalty is not
only an artifact of replaying inputs above five.

At $n=1$ and $n=2$, the original and tree graphs coincide, and the saved
predictions are identical. Small differences in their medians are timing
variation, not evidence of a useful speed or accuracy difference. At larger
$n$, the original AP is not strictly monotonic, and replay does not create
new information. No general claim that increasing slots must improve
detection follows from the table.
